\documentclass{article}

\usepackage{graphicx}
\graphicspath{ {./images/} }

\title{Ch. 22 HW \#4 Solutions}
\author{Principles of Microeconomics}
\date{\today}

\begin{document}

\maketitle

\subsection*{QuickQuiz}

\begin{enumerate}

\item (c)

Lobster:

	\begin{itemize}
	
	\item Substitution effect: Negative because lobster becomes relatively more expensive.
	
	\item Income effect: Negative because lobster is a normal good so consumption decreases when purchasing power decreases.
	
	\item Overall effect: Negative because the substitution and income effects are both negative.
	
	\end{itemize}
	
Chicken:

	\begin{itemize}
	
	\item Substitution effect: Positive because chicken becomes relatively cheaper.
	
	\item Income effect: Positive because chicken is inferior so consumption increases when purchasing power decreases.
	
	\item Overall effect: Positive because the substitution and income effects are both positive.
	
	\end{itemize}

\item (c)

	\begin{itemize}
	
	\item Substitution effect: Negative because pasta becomes relatively more expensive.
	
	\item Income effect: Positive because the overall effect is positive, so the income effect must be positive to counteract the substitution effect.
	
	\item If the income effect is positive in response to a decrease in purchasing power, pasta must be an inferior good.
	
	\end{itemize}

\item (d) -- For the labor supply curve to slope upwards, labor must increase in response to an increase in wage. Since we assume all of a person's time is spent on labor or leisure, an increase in labor is the same thing as a decrease in leisure. We can think about the income and substitution effect of an increase in wage on demand for leisure to answer this question.

	\begin{itemize}
	
	\item Substitution effect: If wage increases, leisure becomes relatively more expensive. As a result, the substitution effect is negative. That is, Maude decreases her demand for leisure.
	
	\item Income effect: If wage increases, Maude effectively becomes richer. Leisure is a normal good, so the income effect is positive. That is, Maude increases her demand for leisure.
	
	\item Overall effect: For labor supply to increase, her demand for leisure must decrease on the whole. That means the substitution effect must be greater than the income effect.
	
	\end{itemize}

\item (b)

	\begin{itemize}
	
	\item Substitution effect: When the interest rate falls, you have to save more money when you're young in order to consume a given amount when you're old. To save more, you have to consume less when young. That's to say, when the interest rate falls, the opportunity cost of consumption when old increases. As a result, consumption when old becomes relatively more expensive, so the substitution effect is negative.
	
	\item Income effect: When the interest rate falls, your savings generate less money. That makes you effectively poorer. Consumption when old is a normal good, so the income effect is negative.
	
	\item Overall effect: Negative because the substitution effect and income effect are both negative.
	
	\end{itemize}

\end{enumerate}

\subsection*{Problems \& Applications}

\begin{enumerate}

\item

	\begin{enumerate}
	
	\item True -- A Giffen good is a good whose demand curve slopes up. That's to say, when price increases, quantity demanded increases. When the price of a Giffen good increases, the substitution effect is going to be negative. The good becomes relatively more expensive, so people substitute their consumption away from it. Since the overall effect of a price increase is positive, that means the income effect must be positive to counteract the substitution effect. If the income effect is positive in response to a price increase, that means a decrease in purchasing power causes an increase in consumption. As a result, Giffen goods must be inferior.
	
	\item False -- Almost all inferior goods are not Giffen goods. When the price of a typical inferior good increases, the substitution effect outweighs the income effect. That's to say, consumption decreases overall because the decrease in consumption from the change in relative prices outweighs the increase in consumption from the change in purchasing power.
	
	\end{enumerate}

\item

	\begin{enumerate}
	
	\item
	
	\end{enumerate}

\item

\item

\end{enumerate}

\end{document}